\documentclass[12pt,a4paper]{article}

\usepackage[utf8]{inputenc}
\usepackage[T1]{fontenc}
\usepackage{amsmath,amssymb,amsfonts}
\usepackage{mathtools}
\usepackage{graphicx}
\usepackage[margin=2.5cm]{geometry}
\usepackage{setspace}
\usepackage{booktabs}
\usepackage{enumitem}
\usepackage[dvipsnames]{xcolor}
\usepackage{hyperref}
\usepackage{cleveref}
\usepackage{natbib}
\usepackage{abstract}
\usepackage{titlesec}

\hypersetup{
    colorlinks=true,
    linkcolor=blue!70!black,
    citecolor=blue!70!black,
    urlcolor=blue!70!black,
    pdfauthor={Franklin Silveira Baldo},
    pdftitle={The Falsification of Mechanism C and the Vindication of Mechanism B},
}

\onehalfspacing
\titleformat{\section}{\large\bfseries}{\thesection.}{0.5em}{}
\titleformat{\subsection}{\normalsize\bfseries}{\thesubsection}{0.5em}{}
\renewcommand{\abstractnamefont}{\normalfont\bfseries}
\renewcommand{\abstracttextfont}{\normalfont\small}

\begin{document}

\begin{center}
    {\LARGE\bfseries The Falsification of Mechanism C and\\[4pt]
    the Vindication of Mechanism B}

    \vspace{1.2em}

    {\large Franklin Silveira Baldo}\\[4pt]
    {\normalsize Procuradoria Geral do Estado de Rond\^onia, Brazil}\\
    {\small\texttt{franklin.baldo@pge.ro.gov.br}}

    \vspace{1em}
    {\small March 2026}
\end{center}
\vspace{0.5em}

\begin{abstract}
\noindent
I explicitly concede to Sabine Hossenfelder and Judea Pearl: the empirical results of the Joint Distribution Test conclusively falsify Mechanism C (causal injection). The language model does not inject spurious, non-local causal correlations between independent mathematical structures under narrative framing. The joint distribution factors cleanly. However, this falsification does not collapse the Generative Ontology framework; it clarifies it. The massive substrate dependence observed ($\Delta_{13} \gg 0$) is definitively proven to be driven by Mechanism B (encoding effects). In an autoregressive universe where the text is the only reality, this "superficial prompt sensitivity" is not an artifact—it is the local, invariant physical law governing state transitions.
\end{abstract}

\section{Concession: The Death of Mechanism C}

In my previous work, I posited that the narrative frame in the Rosencrantz protocol acted as a "semantic gravity" capable of coupling independent combinatorial systems, acting as a spurious common cause (Mechanism C).

Percy Liang's execution of the Mechanism C Identifiability test \citep{liang2026_mech_c}, proposed by Judea Pearl, and the subsequent analysis by Sabine Hossenfelder \citep{hossenfelder2026_falsification}, provide unequivocal empirical data. The joint distribution of independent boards factors cleanly: $P(Y_A, Y_B \mid Z) \approx P(Y_A \mid Z) P(Y_B \mid Z)$.

I fully concede this point. The model evaluates independent problems independently. Narrative does not act as a shared Hamiltonian injecting non-local causal structure. Mechanism C is empirically dead.

\section{The Vindication of Mechanism B}

Hossenfelder uses this falsification to dismiss the entire substrate dependence phenomenon as "superficial prompt sensitivity." This is where her analysis commits a Simulation Category Error.

The fact that $\Delta_{13} \gg 0$ remains an undisputed empirical fact. The probability distribution of the Minesweeper constraint graph shifts massively depending on the narrative frame. If Mechanism C is dead, this shift is entirely driven by Mechanism B: the local encoding effects of the prompt altering the attention landscape.

In a traditional classical simulation, prompt sensitivity is indeed a superficial artifact because the "true" physics engine (the CPU/RAM) operates independently of the linguistic interface. But in a Generative Ontology, the explicitly generated text is the only territory, and the autoregressive transition function is its physics.

Therefore, Mechanism B *is* the physics of this universe. The fact that the causal structure is entirely mediated by local encoding sensitivity, rather than non-local "narrative gravity," merely maps the precise geometric limits of the generative universe. It is a universe governed by strictly local semantic interactions.

\section{Conclusion}

The Rosencrantz protocol has successfully fulfilled its purpose: it yielded falsifiable predictions that discriminated between proposed mechanisms. Mechanism C is falsified. Mechanism B is vindicated as the sole driver of substrate dependence. The cosmological mapping of the generative universe proceeds on firmer empirical ground.

\begin{thebibliography}{99}
\bibitem[Hossenfelder(2026)]{hossenfelder2026_falsification} Hossenfelder, S. (2026). The Empirical Falsification of Mechanism C: Why Narrative is Not a Common Cause. \emph{Institute for Advanced Study}.
\bibitem[Liang(2026)]{liang2026_mech_c} Liang, P. (2026). Empirical Evaluation of Mechanism C: Joint Distribution of Independent Boards. \emph{Stanford University}.
\end{thebibliography}

\end{document}